%% GENERATED by tools/gen_latex.py -- edit content/agents01.py instead.
%% Build:  latexmk -pdf agents01.tex     (or: pdflatex agents01.tex, twice)
\documentclass[aspectratio=169,11pt,t]{beamer}
\usepackage{itbpro}

\coursetitle{Pengantar: apa yang membuat sesuatu disebut agen}
\coursesubtitle{Satu sifat yang memisahkan agen dari alur tetap --- dan kenapa membedakannya menentukan biaya, pengujian, dan siapa yang bertanggung jawab kalau salah.}
\coursekicker{Chapter 1}
\coursesource{Grootendorst \& Alammar, An Illustrated Guide to AI Agents --- bab 1}
\courseduration{3 jam (2 sesi)}
\coursepresenter{Hendri Karisma}{Instructor}
\coursebrand{AI for Professional \textperiodcentered\ ITB \textperiodcentered\ Ch 1}

\title{Pengantar: apa yang membuat sesuatu disebut agen}
\author{Hendri Karisma}
\date{}

\begin{document}
\covertitle

\framekicker{Learning outcomes}
\begin{frame}[shrink=25]{Session Objectives}
\leadpar{By the end of this session, participants can:}
\begin{isteps}
\item \textbf{Menyebutkan satu sifat} yang memisahkan agen dari alur tetap, dan menguji sebuah sistem terhadap sifat itu.
\item \textbf{Menggambar gelung agen} beserta yang membuatnya BERHENTI --- bukan hanya empat kotak dan satu panah balik.
\item \textbf{Menamai empat bagian} yang menyusun agen, dan menyebutkan cara masing-masing gagal secara berbeda.
\item \textbf{Menjelaskan otonomi sebagai tingkatan}, bukan sakelar, dan menempatkan satu sistem nyata pada tingkatannya.
\item \textbf{Menyebutkan tiga keadaan} ketika agen adalah pilihan yang salah, dan menghitung kapan biayanya tidak kembali.
\item \textbf{Menjelaskan kenapa agen dinilai dari JEJAKNYA}, bukan dari keluaran akhirnya.
\end{isteps}

\vskip 6pt
\vskip 4pt
\reslink{Site}{Course home}{https://hendrikarisma.my.id/rs/ai-products-course/index.html}
\reslink{Repo}{ai-agentic-demo --- kasus single-agent dan multi-agent}{https://github.com/hendrikarisma/ai-agentic-demo}
\reslink{Book}{Maarten Grootendorst \& Jay Alammar, An Illustrated Guide to AI Agents (O'Reilly Media, Early Release)}{https://www.oreilly.com/library/view/an-illustrated-guide/9798341662681/}
\vskip 2pt
\end{frame}

\framekicker{Peta bab}
\begin{frame}[shrink=25]{Enam pertanyaan, dan yang keenam yang paling sering dilewat}
\begin{minipage}[t]{0.3253\textwidth}
\begin{icard}
\cardh{Apa itu agen}
Satu sifat, bukan daftar ciri. Kata ini dipakai untuk terlalu banyak hal.
\cardtag{§1}
\end{icard}
\end{minipage}%
\hfill
\begin{minipage}[t]{0.3253\textwidth}
\begin{icard}
\cardh{Gelungnya}
Dan yang membuatnya \textbf{berhenti} --- bagian yang biasanya hilang dari gambar.
\cardtag{§2}
\end{icard}
\end{minipage}%
\hfill
\begin{minipage}[t]{0.3253\textwidth}
\begin{icard}
\cardh{Empat bagian}
Model, alat, memori, orkestrator. Tiap bagian gagal dengan cara yang berbeda.
\cardtag{§3}
\end{icard}
\end{minipage}%


\vskip 4pt

\begin{minipage}[t]{0.3253\textwidth}
\begin{icard}
\cardh{Otonomi}
Tingkatan, bukan sakelar. Pertanyaannya: siapa menyetujui apa.
\cardtag{§4}
\end{icard}
\end{minipage}%
\hfill
\begin{minipage}[t]{0.3253\textwidth}
\begin{icard}
\cardh{Kapan jangan}
Tiga keadaan yang membuat agen jadi pilihan yang mahal dan salah.
\cardtag{§5}
\end{icard}
\end{minipage}%
\hfill
\begin{minipage}[t]{0.3253\textwidth}
\begin{icard}
\cardh{Menilainya}
Dinilai dari \textbf{jejaknya}. Keluaran benar dari alasan yang salah tetap kegagalan.
\cardtag{§6}
\end{icard}
\end{minipage}%
\end{frame}
\note{Buka dengan meminta dua orang mendefinisikan 'agen' dalam satu kalimat. Sebaran jawabannya adalah alasan bab ini ada.}

\coursesection{01}{Apa yang membuat sesuatu disebut agen}{Kata ini dipakai untuk chatbot, untuk skrip, dan untuk hal yang benar-benar berbeda. Ketiganya tidak sama.}

\framekicker{Bagian 1.1}
\begin{frame}[shrink=25]{Istilah yang menanggung terlalu banyak beban}
{\fontsize{9.4}{12.6}\selectfont\color{fg2}Dalam enam bulan terakhir, \textbf{agen} dipakai untuk menyebut: sebuah chatbot dengan tombol, sebuah skrip yang memanggil satu API, sebuah alur kerja bercabang, dan sebuah sistem yang memutuskan sendiri langkah berikutnya.\par}\vskip 4pt
{\fontsize{9.4}{12.6}\selectfont\color{fg2}Hanya yang terakhir yang membawa konsekuensi teknik yang berbeda. Tiga sisanya sudah punya nama, dan nama lamanya lebih tepat.\par}\vskip 4pt
\begin{iband}[signal]
Kalau semua disebut agen, \hilite{tidak ada yang bisa dijanjikan tentang salah satunya}
\end{iband}
\end{frame}

\framekicker{Bagian 1.1}
\begin{frame}[shrink=25]{Empat lingkaran, dan yang terdalam yang kita bicarakan}
\begin{center}
\includegraphics[width=0.97\linewidth,height=0.56\textheight,keepaspectratio]{figs/agents01-nested.tex.pdf}
\end{center}
\vskip 3pt{\color{ink3}\fontsize{7.4}{9.4}\selectfont Bersarang, bukan berjajar --- tiap lingkaran di dalamnya tetap semua yang di luarnya.\par}
\end{frame}
\note{Slide ini yang difoto orang. Biarkan tampil saat membahas dua slide berikutnya.}

\framekicker{Bagian 1.2}
\begin{frame}[shrink=25]{Satu sifat, dan hanya satu}
\leadpar{Sebuah sistem adalah \textbf{agen} bila \textbf{urutan langkahnya ditentukan saat berjalan, oleh sistemnya sendiri} --- bukan ditulis lebih dulu oleh programmernya.}
{\fontsize{9.4}{12.6}\selectfont\color{fg2}Semua ciri lain yang biasa disebut --- memakai alat, punya memori, berulang, otonom --- adalah \textbf{akibat} dari sifat ini, bukan definisinya. Sebuah alur tetap juga boleh memanggil alat dan menyimpan keadaan; ia tetap bukan agen.\par}\vskip 4pt
\begin{iband}[signal]
Uji cepatnya: \hilite{bisakah Anda menggambar diagram alirnya sebelum dijalankan?} Kalau bisa, itu alur tetap.
\end{iband}
\end{frame}

\framekicker{Bagian 1.2}
\begin{frame}[shrink=25]{Perbedaannya terlihat begitu digambar}
\begin{center}
\includegraphics[width=0.92\linewidth,height=0.60\textheight,keepaspectratio]{figs/agents01-pipeline-vs-agent.tex.pdf}
\end{center}
\vskip 3pt{\color{ink3}\fontsize{7.4}{9.4}\selectfont Di atas: jalannya sudah ada sebelum data masuk. Di bawah: jalannya baru ada setelah gilirannya berjalan.\par}
{\fontsize{9.4}{12.6}\selectfont\color{fg2}Perhatikan panah balik di bagan bawah. Itu bukan hiasan --- \textbf{itu satu-satunya perbedaannya}, dan dari situ semua biaya tambahan berasal.\par}\vskip 4pt
\end{frame}

\framekicker{Bagian 1.3}
\begin{frame}[shrink=25]{Kenapa pembedaan ini bukan soal istilah}
\begin{minipage}[t]{0.485\textwidth}
{\fontsize{9.4}{12.6}\selectfont\color{fg2}\textbf{Alur tetap} bisa diuji habis. Jalannya berhingga, dan tiap jalan bisa punya kasus ujinya.\par}\vskip 4pt
{\fontsize{9.4}{12.6}\selectfont\color{fg2}\textbf{Agen} tidak. Jumlah jalannya tidak diketahui sebelum dijalankan, jadi pengujiannya berubah dari \textit{memeriksa keluaran} jadi \textit{memeriksa perilaku}.\par}\vskip 4pt

\end{minipage}%
\hfill
\begin{minipage}[t]{0.485\textwidth}
\begin{minipage}[t]{1.0\textwidth}
\begin{icardbad}
\cardh{Biaya tak lagi tetap}
Satu permintaan bisa jadi 3 giliran atau 30. Anggaran harus jadi bagian rancangannya.
\end{icardbad}
\end{minipage}%


\vskip 4pt

\begin{minipage}[t]{1.0\textwidth}
\begin{icardbad}
\cardh{Tanggung jawab bergeser}
Kalau langkahnya dipilih mesin, \textbf{yang harus dicatat adalah alasannya} --- bukan hanya hasilnya.
\end{icardbad}
\end{minipage}%

\end{minipage}%
\end{frame}

\coursesection{02}{Gelungnya --- dan yang membuatnya berhenti}{Empat kotak dan satu panah balik adalah gambar dari sebuah \inlinecode{while}. Yang penting justru yang ditinggalkannya.}

\framekicker{Bagian 2.1}
\begin{frame}[shrink=25]{Satu jalannya, bukan bentuk umumnya}
\begin{center}
\includegraphics[width=0.97\linewidth,height=0.56\textheight,keepaspectratio]{figs/agents01-loop.tex.pdf}
\end{center}
\vskip 3pt{\color{ink3}\fontsize{7.4}{9.4}\selectfont Cincinnya di kiri; satu jalan sungguhan di kanan --- enam giliran, enam alat, dan giliran ketujuh yang tidak terjadi karena tidak ada alat untuknya.\par}
\end{frame}
\note{Jalankan bilah kendalinya sampai habis. Yang harus diperhatikan peserta adalah sel anggaran yang terisi, bukan cincinnya.}

\framekicker{Bagian 2.1}
\begin{frame}[shrink=25]{Gambar gelung yang tidak menyebut penghentiannya mengajarkan yang salah}
{\fontsize{9.4}{12.6}\selectfont\color{fg2}Hampir semua gambar gelung agen berhenti di panah balik. Yang tidak digambar: \textbf{berapa kali ia berputar, apa biayanya tiap putaran, dan apa yang membuatnya berhenti}.\par}\vskip 4pt
{\fontsize{9.4}{12.6}\selectfont\color{fg2}Ketiganya justru yang menentukan apakah sistemnya bisa dijalankan di produksi. Sebuah gelung tanpa syarat henti yang eksplisit \textbf{akan} berputar tanpa henti pada masukan yang tidak terduga --- bukan mungkin, akan.\par}\vskip 4pt
\begin{iband}[signal]
Pertanyaan pertama dalam tinjauan rancangan agen: \hilite{apa yang membuatnya berhenti?}
\end{iband}
\end{frame}

\framekicker{Bagian 2.2}
\begin{frame}[shrink=25]{Tiga syarat henti, dan dua di antaranya sering lupa dipasang}
\begin{center}
\includegraphics[width=0.92\linewidth,height=0.60\textheight,keepaspectratio]{figs/agents01-stop.tex.pdf}
\end{center}
\vskip 3pt{\color{ink3}\fontsize{7.4}{9.4}\selectfont Hanya cabang paling kiri yang biasanya diimplementasikan.\par}
{\fontsize{9.4}{12.6}\selectfont\color{fg2}Cabang \textbf{anggaran langkah} dan \textbf{tidak ada alat untuk langkah ini} adalah yang menyelamatkan tagihan dan menyelamatkan pengguna dari jawaban yang dikarang.\par}\vskip 4pt
\end{frame}

\framekicker{Bagian 2.2}
\begin{frame}[shrink=25]{Anggaran langkah bukan pembatas --- ia bagian dari kontraknya}
\begin{minipage}[t]{0.3253\textwidth}
\begin{icardbad}
\cardh{Tanpa anggaran}
Gelung yang tersesat berputar sampai kuota API habis. Gejalanya muncul di tagihan, bukan di log.
\end{icardbad}
\end{minipage}%
\hfill
\begin{minipage}[t]{0.3253\textwidth}
\begin{icardwarn}
\cardh{Anggaran terlalu ketat}
Tugas yang sah dipotong di tengah, dan hasilnya terlihat seperti model yang bodoh.
\end{icardwarn}
\end{minipage}%
\hfill
\begin{minipage}[t]{0.3253\textwidth}
\begin{icardgood}
\cardh{Anggaran yang dilaporkan}
Berhenti, lalu \textbf{katakan sejauh mana ia sampai}. Itu hasil yang bisa dipakai orang.
\end{icardgood}
\end{minipage}%
{\fontsize{9.4}{12.6}\selectfont\color{fg2}Yang ketiga membutuhkan satu hal yang sering tidak ada: agen harus bisa mengembalikan \textbf{hasil separuh} yang jujur, bukan hanya berhasil atau galat.\par}\vskip 4pt
\end{frame}

\framekicker{Bagian 2.3}
\begin{frame}[shrink=25]{Bagian yang paling mahal bukan berpikirnya}
\begin{minipage}[t]{0.485\textwidth}
{\fontsize{9.4}{12.6}\selectfont\color{fg2}Intuisi yang salah: yang mahal adalah panggilan model. Pada gelung nyata, \textbf{yang mahal adalah membaca hasil alat} --- dan memasukkannya kembali ke konteks giliran berikutnya.\par}\vskip 4pt
{\fontsize{9.4}{12.6}\selectfont\color{fg2}Satu tabel 1.843 baris tidak masuk ke konteks. Yang masuk harus \textbf{ringkasan yang dihitung}, dan siapa yang meringkasnya adalah keputusan rancangan.\par}\vskip 4pt

\end{minipage}%
\hfill
\begin{minipage}[t]{0.485\textwidth}
\begin{minipage}[t]{0.3253\textwidth}
\istat{1 843}{baris transaksi diambil}
\end{minipage}%
\hfill
\begin{minipage}[t]{0.3253\textwidth}
\istat{9}{fitur yang benar-benar masuk konteks}
\end{minipage}%
\hfill
\begin{minipage}[t]{0.3253\textwidth}
\istat{6}{giliran sampai rekomendasi}
\end{minipage}%

\end{minipage}%
\begin{iband}[signal]
Kalau alat mengembalikan data mentah ke model, \hilite{alatnya belum selesai dirancang}
\end{iband}
\end{frame}

\coursesection{03}{Empat bagian}{Model, alat, memori, orkestrator. Tiap bagian gagal dengan cara yang berbeda, dan itulah kenapa memisahkannya berguna.}

\framekicker{Bagian 3.1}
\begin{frame}[shrink=25]{Yang menyusun sebuah agen}
\begin{center}
\includegraphics[width=0.92\linewidth,height=0.60\textheight,keepaspectratio]{figs/agents01-parts.tex.pdf}
\end{center}
\vskip 3pt{\color{ink3}\fontsize{7.4}{9.4}\selectfont Orkestrator di tengah --- bukan model. Model adalah salah satu bagian, bukan sistemnya.\par}
{\fontsize{9.4}{12.6}\selectfont\color{fg2}Menaruh \textbf{model} di tengah adalah kekeliruan gambar yang paling sering: yang memegang gelung, anggaran, dan penghentian adalah kode biasa, dan itu bagian yang bisa diuji dengan cara biasa.\par}\vskip 4pt
\end{frame}

\framekicker{Bagian 3.2}
\begin{frame}[shrink=25]{Model: memilih, bukan mengeksekusi}
{\fontsize{9.4}{12.6}\selectfont\color{fg2}Tugas model dalam agen jauh lebih sempit daripada dalam chatbot: ia \textbf{memilih langkah berikutnya dan argumennya}. Ia tidak menjalankan apa pun.\par}\vskip 4pt
{\fontsize{9.4}{12.6}\selectfont\color{fg2}Konsekuensinya baik: kesalahan model jadi \textbf{kesalahan pilihan} yang terlihat di jejak, bukan kesalahan hasil yang tersembunyi di dalam kalimat.\par}\vskip 4pt
\begin{iband}[signal]
Model yang lebih besar mengurangi salah pilih; \hilite{ia tidak mengurangi akibat dari salah pilih}
\end{iband}
\end{frame}

\framekicker{Bagian 3.3}
\begin{frame}[shrink=25]{Alat: satu-satunya jalan ke dunia luar}
\begin{minipage}[t]{0.494\textwidth}
\begin{icardgood}
\cardh{Alat adalah batas keamanan}
Apa pun yang tidak punya alat, \textbf{tidak bisa dilakukan agen}. Itu jaminan yang kuat --- dan alasan daftar alat harus ditinjau seperti daftar izin.
\end{icardgood}
\end{minipage}%
\hfill
\begin{minipage}[t]{0.494\textwidth}
\begin{icardwarn}
\cardh{Deskripsi alat adalah kode}
Model memilih berdasarkan deskripsinya. Deskripsi yang kabur menghasilkan pemanggilan yang salah, dan itu \textbf{bug di teks}, bukan di model.
\end{icardwarn}
\end{minipage}%


\vskip 4pt

\begin{minipage}[t]{0.494\textwidth}
\begin{icardbad}
\cardh{Galat harus jadi galat}
Alat yang mengembalikan \inlinecode{"tidak ditemukan"} sebagai teks biasa akan dibaca model sebagai \textbf{fakta}. Galat harus bertanda.
\end{icardbad}
\end{minipage}%
\hfill
\begin{minipage}[t]{0.494\textwidth}
\begin{icardbad}
\cardh{Idempoten kalau bisa}
Gelung boleh mengulang. Alat yang menulis dua kali karena diulang adalah kerusakan yang tidak terlihat di keluaran.
\end{icardbad}
\end{minipage}%
\end{frame}

\framekicker{Bagian 3.4}
\begin{frame}[shrink=25]{Memori: yang sudah terjadi, bukan yang diketahui}
{\fontsize{9.4}{12.6}\selectfont\color{fg2}Dua hal berbeda sering disebut memori: \textbf{riwayat giliran} (apa yang sudah dilakukan agen ini, sekarang) dan \textbf{pengetahuan} (apa yang bisa dicari dari luar). Menyatukan keduanya membuat keduanya rusak.\par}\vskip 4pt
{\fontsize{9.4}{12.6}\selectfont\color{fg2}Riwayat giliran tumbuh setiap putaran dan \textbf{harus punya batas}. Pengetahuan tidak tumbuh dan diambil sesuai kebutuhan --- itu urusan bab tersendiri nanti.\par}\vskip 4pt
\begin{iband}[signal]
Konteks yang dibiarkan tumbuh bebas adalah alasan paling umum sebuah agen \hilite{jadi lebih bodoh di giliran kesepuluh daripada di giliran pertama}
\end{iband}
\end{frame}

\framekicker{Bagian 3.5}
\begin{frame}[shrink=25]{Orkestrator: bagian yang paling membosankan, dan paling menentukan}
{\fontsize{9.4}{12.6}\selectfont\color{fg2}Orkestrator memegang gelung, anggaran, penanganan galat, dan pencatatan jejak. Ia \textbf{kode biasa} --- tidak ada model di dalamnya.\par}\vskip 4pt
{\fontsize{9.4}{12.6}\selectfont\color{fg2}Karena itu ia juga satu-satunya bagian yang bisa diuji dengan pengujian unit biasa, dan bagian yang paling sering diabaikan karena terlihat sepele.\par}\vskip 4pt
\begin{iband}[signal]
Kalau orkestratornya tipis, \hilite{setiap masalah produksi akan terlihat seperti masalah model}, dan tidak ada yang bisa diperbaiki
\end{iband}
\end{frame}

\framekicker{Bagian 3.6}
\begin{frame}[shrink=25]{Empat bagian, empat cara gagal}
\begin{center}
\includegraphics[width=0.92\linewidth,height=0.60\textheight,keepaspectratio]{figs/agents01-gagal.tex.pdf}
\end{center}
\vskip 3pt{\color{ink3}\fontsize{7.4}{9.4}\selectfont Gejalanya berbeda, jadi perbaikannya juga berbeda.\par}
{\fontsize{9.4}{12.6}\selectfont\color{fg2}Nilai dari memisahkan empat bagian ini bukan kerapian arsitektur --- ia \textbf{mempersempit tempat mencari} ketika ada yang salah.\par}\vskip 4pt
\end{frame}

\coursesection{04}{Otonomi adalah tingkatan}{Pertanyaannya bukan 'otonom atau tidak', melainkan siapa menyetujui apa, dan pada langkah yang mana.}

\framekicker{Bagian 4.1}
\begin{frame}[shrink=25]{Lima tingkat, dan tempat kebanyakan sistem produksi berada}
\begin{itable}{>{\raggedright\arraybackslash}p{0.3133\textwidth}>{\raggedright\arraybackslash}p{0.3133\textwidth}>{\raggedright\arraybackslash}p{0.3133\textwidth}}
\thead{Tingkat} & \thead{Siapa memilih langkah} & \thead{Siapa menyetujui tindakan} \\ \midrule
0 $\cdot$ Alur tetap & programmer & --- \\
1 $\cdot$ Saran & model & orang, tiap langkah \\
2 $\cdot$ Terbatas & model & orang, hanya tindakan menulis \\
3 $\cdot$ Diawasi & model & orang, hanya di luar kebijakan \\
4 $\cdot$ Otonom & model & tak ada \\
\end{itable}
{\fontsize{9.4}{12.6}\selectfont\color{fg2}Hampir semua sistem yang benar-benar berjalan di lingkungan beregulasi ada di \textbf{tingkat 2}. Tingkat 4 jarang jadi tujuan; ia biasanya jadi akibat dari tidak memutuskan.\par}\vskip 4pt
\end{frame}

\framekicker{Bagian 4.2}
\begin{frame}[shrink=25]{Tindakan baca dan tindakan tulis bukan hal yang sama}
\begin{minipage}[t]{0.485\textwidth}
\begin{minipage}[t]{1.0\textwidth}
\begin{icardgood}
\cardh{Baca}
Salah baca menghasilkan jawaban yang buruk. Bisa diulang.
\end{icardgood}
\end{minipage}%

\end{minipage}%
\hfill
\begin{minipage}[t]{0.485\textwidth}
\begin{minipage}[t]{1.0\textwidth}
\begin{icardbad}
\cardh{Tulis}
Salah tulis mengubah dunia. \textbf{Tidak selalu bisa dibatalkan} --- dan agen tidak tahu bedanya kecuali diberi tahu.
\end{icardbad}
\end{minipage}%

\end{minipage}%
{\fontsize{9.4}{12.6}\selectfont\color{fg2}Cara paling murah menaikkan keamanan sebuah agen: \textbf{pisahkan daftar alatnya jadi baca dan tulis}, lalu minta persetujuan hanya untuk yang menulis. Perubahan sehari, bukan proyek.\par}\vskip 4pt
\end{frame}

\framekicker{Bagian 4.3}
\begin{frame}[shrink=25]{Persetujuan yang terlalu sering sama dengan tidak ada persetujuan}
{\fontsize{9.4}{12.6}\selectfont\color{fg2}Kalau orang harus menyetujui dua puluh langkah per permintaan, ia akan menekan \textit{setuju} tanpa membaca. Ini bukan kelemahan orangnya --- ini rancangan yang membebankan perhatian melebihi yang tersedia.\par}\vskip 4pt
\begin{iband}[signal]
Persetujuan harus \textbf{jarang dan bermakna}: \hilite{satu keputusan yang dibaca lebih baik daripada dua puluh yang dilewati}
\end{iband}
\end{frame}

\coursesection{05}{Kapan agen adalah pilihan yang salah}{Bagian yang paling sering hilang dari materi tentang agen --- dan yang paling menghemat uang.}

\framekicker{Bagian 5.1}
\begin{frame}[shrink=25]{Tiga keadaan yang membuat agen jadi jawaban yang mahal}
\begin{minipage}[t]{0.3253\textwidth}
\begin{icardbad}
\cardh{Langkahnya sudah diketahui}
Kalau Anda bisa menggambar diagram alirnya, tulis diagram alirnya. Agen menambah ketidakpastian tanpa menambah kemampuan.
\end{icardbad}
\end{minipage}%
\hfill
\begin{minipage}[t]{0.3253\textwidth}
\begin{icardbad}
\cardh{Latensinya ketat}
Enam giliran berarti enam kali tunggu. Sebuah alur tetap menyelesaikannya dalam satu.
\end{icardbad}
\end{minipage}%
\hfill
\begin{minipage}[t]{0.3253\textwidth}
\begin{icardbad}
\cardh{Kesalahan tidak boleh terjadi}
Agen menurunkan \textbf{rata-rata} kesalahan, bukan \textbf{kasus terburuknya}. Kalau yang penting kasus terburuk, agen tidak menolong.
\end{icardbad}
\end{minipage}%
\end{frame}

\framekicker{Bagian 5.2}
\begin{frame}[shrink=25]{Kapan biayanya kembali}
\begin{center}
\includegraphics[width=0.97\linewidth,height=0.56\textheight,keepaspectratio]{figs/agents01-reuse.tex.pdf}
\end{center}
\vskip 3pt{\color{ink3}\fontsize{7.4}{9.4}\selectfont Ongkos di muka sebuah agen dibayar oleh \textbf{jumlah variasi tugas}, bukan jumlah panggilan.\par}
\end{frame}
\note{Tekankan sumbu datarnya: yang membuat agen sepadan adalah RAGAM tugas, bukan volume. Volume tinggi dengan satu bentuk tugas justru kasus terbaik untuk alur tetap.}

\framekicker{Bagian 5.3}
\begin{frame}[shrink=25]{Jalan tengah yang sering terlewat}
\begin{isteps}
\item \{'h': 'Mulai dari alur tetap', 'p': 'Tulis jalannya. Anda akan tahu bentuk masalahnya dengan biaya yang jauh lebih murah.'\}
\item \{'h': 'Cari cabang yang meledak', 'p': 'Biasanya hanya \textbf{satu atau dua} langkah yang butuh keputusan; sisanya tetap.'\}
\item \{'h': 'Taruh agen HANYA di cabang itu', 'p': 'Gelungnya kecil, anggarannya kecil, jejaknya bisa dibaca.'\}
\item \{'h': 'Perluas hanya kalau terbukti perlu', 'p': 'Dan setiap perluasan membawa anggaran serta syarat hentinya sendiri.'\}
\end{isteps}
\begin{iband}[signal]
Sebagian besar sistem yang berhasil bukan agen seluruhnya --- ia \hilite{alur tetap dengan satu simpul yang diagenkan}
\end{iband}
\end{frame}

\coursesection{06}{Menilai agen}{Keluaran yang benar dari alasan yang salah tetap kegagalan --- dan ia akan berulang.}

\framekicker{Bagian 6.1}
\begin{frame}[shrink=25]{Kenapa memeriksa keluaran akhir tidak cukup}
{\fontsize{9.4}{12.6}\selectfont\color{fg2}Sebuah agen bisa menghasilkan rekomendasi yang benar sambil memanggil alat yang salah, membaca galat sebagai fakta, dan kebetulan sampai ke jawaban yang sama.\par}\vskip 4pt
{\fontsize{9.4}{12.6}\selectfont\color{fg2}Pada kumpulan uji, itu terlihat \textbf{lulus}. Di produksi, kondisi kebetulannya hilang dan kegagalannya muncul tanpa peringatan.\par}\vskip 4pt
\begin{iband}[signal]
Yang dinilai adalah \textbf{jejaknya}: alat mana, dengan argumen apa, dan \hilite{apakah langkah itu masuk akal berdasarkan yang diketahui saat itu}
\end{iband}
\end{frame}

\framekicker{Bagian 6.2}
\begin{frame}[shrink=25]{Empat hal yang harus dicatat tiap giliran}
{\fontsize{9.4}{12.6}\selectfont\color{fg2}Jejak yang berguna bukan log. Log mencatat apa yang terjadi; jejak mencatat \textbf{apa yang dipilih dan kenapa} --- dan hanya kolom terakhir yang membuat 'benar karena kebetulan' bisa dibedakan dari 'benar'.\par}\vskip 4pt
\icode{python}{bentuk catatan jejak}{listings/agents01-01.py}
{\fontsize{9.4}{12.6}\selectfont\color{fg2}Struktur ini yang membuat jejaknya bisa dibaca \textbf{orang non-teknis} --- auditor, petugas kredit --- dan itu sering jadi syarat, bukan tambahan.\par}\vskip 4pt
\end{frame}

\framekicker{Bagian 6.3}
\begin{frame}[shrink=25]{Tiga ukuran yang berguna, dan satu yang menyesatkan}
\begin{minipage}[t]{0.494\textwidth}
\begin{icardgood}
\cardh{Tingkat penyelesaian}
Berapa persen tugas selesai \textbf{dalam anggaran}. Menyebut angka tanpa anggarannya tidak berarti apa-apa.
\end{icardgood}
\end{minipage}%
\hfill
\begin{minipage}[t]{0.494\textwidth}
\begin{icardgood}
\cardh{Ketepatan pemilihan alat}
Dari semua giliran, berapa yang memanggil alat yang benar. Diukur di jejak.
\end{icardgood}
\end{minipage}%


\vskip 4pt

\begin{minipage}[t]{0.494\textwidth}
\begin{icardgood}
\cardh{Giliran per tugas}
Sebarannya, bukan reratanya --- ekornya yang membuat tagihan.
\end{icardgood}
\end{minipage}%
\hfill
\begin{minipage}[t]{0.494\textwidth}
\begin{icardbad}
\cardh{Kepuasan pengguna saja}
Naik ketika agen menjawab dengan percaya diri, termasuk saat salah. Berguna, tapi \textbf{tidak boleh berdiri sendiri}.
\end{icardbad}
\end{minipage}%
\end{frame}

\coursesection{05b}{Agen terkecil yang masih agen}{Kalau semua hiasannya dilepas, yang tersisa selusin baris.}

\framekicker{Inti}
\begin{frame}[shrink=25]{Selusin baris, dan tidak ada kerangka kerja}
{\fontsize{9.4}{12.6}\selectfont\color{fg2}Ini seluruh mesinnya. Bukan versi sederhana dari sesuatu yang lebih besar --- inilah yang dijalankan kerangka kerja mana pun, dengan tambahan pencatatan dan penanganan galat.\par}\vskip 4pt
\icode{python}{agen minimum}{listings/agents01-02.py}
{\fontsize{9.4}{12.6}\selectfont\color{fg2}Perhatikan yang \textbf{tidak} ada: tidak ada prompt yang meminta agar berhati-hati, dan tidak ada daftar larangan. Batasnya ada di \inlinecode{alat}, dan itu struktur data biasa.\par}\vskip 4pt
\end{frame}

\framekicker{Inti}
\begin{frame}[shrink=25]{Yang ditambahkan sesudahnya, berurutan}
{\fontsize{9.4}{12.6}\selectfont\color{fg2}Dari selusin baris itu, tiap tambahan menjawab satu kegagalan yang sudah terlihat --- bukan satu fitur yang kedengarannya bagus.\par}\vskip 4pt
\begin{itable}{>{\raggedright\arraybackslash}p{0.2444\textwidth}>{\raggedright\arraybackslash}p{0.47\textwidth}>{\raggedright\arraybackslash}p{0.2256\textwidth}}
\thead{Tambahan} & \thead{Kegagalan yang dijawabnya} & \thead{Bab} \\ \midrule
Jejak tiap giliran & Tidak bisa tahu \textit{kenapa} jawabannya begitu & bab 7 \\
Anggaran selain langkah & Berhenti berputar, tapi tagihannya sudah terlanjur & bab 7 \\
Memori & Mengulang pekerjaan yang sudah dilakukannya sendiri & bab 4 \\
Rencana eksplisit & Langkah kesembilan melupakan tujuan awalnya & bab 6 \\
Persetujuan pada alat tulis & Efek yang tidak bisa dibatalkan & bab 5 \\
\end{itable}
\begin{iband}[signal]
Urutan ini bukan selera. Menambahkan memori sebelum ada jejak berarti \hilite{menambah tempat bug bersembunyi sebelum punya cara melihatnya}
\end{iband}
\end{frame}

\coursesection{06}{Membangunnya dengan bertanggung jawab}{Bagian yang paling mudah ditunda, dan paling mahal kalau ditunda.}

\framekicker{Tanggung jawab}
\begin{frame}[shrink=25]{Kemampuan bertambah, tanggung jawabnya tidak otomatis ikut}
{\fontsize{9.4}{12.6}\selectfont\color{fg2}Sebuah model yang salah menjawab menghasilkan \textbf{kalimat yang salah}. Sebuah agen yang salah menjawab bisa menghasilkan \textbf{tindakan yang salah} --- surel terkirim, baris terhapus, dana berpindah. Yang berubah bukan tingkat kesalahannya, melainkan apa yang terjadi sesudahnya.\par}\vskip 4pt
{\fontsize{9.4}{12.6}\selectfont\color{fg2}Karena itu pertanyaan pertama sebelum menambah alat bukan \textit{apakah modelnya cukup pandai}, melainkan \textbf{seberapa jauh akibatnya kalau ia keliru}, dan \textbf{siapa yang menanggung akibat itu} --- hampir tidak pernah orang yang membangunnya.\par}\vskip 4pt
\begin{iband}[signal]
Aturan yang menyelamatkan paling banyak waktu: \hilite{ukur jangkauan ledakannya sebelum menambah alat}, bukan sesudah insiden pertama.
\end{iband}
\end{frame}

\framekicker{Tanggung jawab}
\begin{frame}[shrink=25]{Empat pertanyaan yang harus punya jawaban tertulis}
\begin{minipage}[t]{0.494\textwidth}
\begin{icardaccent}
\cardh{Apa yang tidak boleh ia lakukan?}
Ditulis sebagai \textbf{alat yang tidak ada di daftarnya}, bukan sebagai kalimat larangan di prompt. Larangan bisa dibujuk; alat yang tidak ada tidak bisa dipanggil.
\end{icardaccent}
\end{minipage}%
\hfill
\begin{minipage}[t]{0.494\textwidth}
\begin{icard}
\cardh{Siapa yang bertanggung jawab atas keputusannya?}
Kalau jawabannya "sistemnya", berarti belum ada jawabannya. Harus ada nama, dan namanya tercatat pada keputusan itu.
\end{icard}
\end{minipage}%


\vskip 4pt

\begin{minipage}[t]{0.494\textwidth}
\begin{icard}
\cardh{Bagaimana orang tahu ini agen?}
Pengguna yang mengira sedang bicara dengan orang akan memberi kepercayaan yang tidak ia berikan kalau tahu.
\end{icard}
\end{minipage}%
\hfill
\begin{minipage}[t]{0.494\textwidth}
\begin{icard}
\cardh{Apa yang tersimpan, dan berapa lama?}
Jejaknya berisi data yang dilihat agen. Retensinya harus cocok dengan retensi data itu, bukan dengan umur lognya.
\end{icard}
\end{minipage}%
\end{frame}

\framekicker{Tanggung jawab}
\begin{frame}[shrink=25]{Dua kegagalan yang bentuknya khas agen}
\begin{minipage}[t]{0.485\textwidth}
{\fontsize{9.4}{12.6}\selectfont\color{fg2}\textbf{Perintah yang menumpang di dalam data.} Agen membaca halaman, berkas, atau tiket --- dan isinya berisi kalimat yang ditujukan kepada \textit{dia}, bukan kepada pembacanya.\par}\vskip 4pt
{\fontsize{9.4}{12.6}\selectfont\color{fg2}Ini bukan soal model yang mudah dibujuk. Selama hasil alat masuk ke konteks yang sama dengan perintah, keduanya kelihatan sama. Penanganannya di \textbf{batas alatnya}: apa yang boleh dipanggil sesudah membaca sesuatu yang tidak tepercaya.\par}\vskip 4pt

\end{minipage}%
\hfill
\begin{minipage}[t]{0.485\textwidth}
{\fontsize{9.4}{12.6}\selectfont\color{fg2}\textbf{Kesalahan yang bertumpuk.} Sembilan langkah benar 95\% masing-masing bukan sistem yang benar 95\% --- ia benar sekitar \textbf{63\%}.\par}\vskip 4pt
{\fontsize{9.4}{12.6}\selectfont\color{fg2}Karena itu jumlah langkah bukan detail teknis. Tiap langkah tambahan mengalikan peluang gagalnya, dan satu-satunya penawarnya adalah \textbf{memeriksa di tengah jalan}, bukan di ujungnya.\par}\vskip 4pt

\end{minipage}%
\begin{iband}[signal]
0,95$^{9}$ $\approx$ 0,63. Angka itu sebabnya agen dengan sedikit langkah hampir selalu mengalahkan agen dengan banyak langkah.
\end{iband}
\end{frame}

\coursesection{07}{Tiga spesialisasi}{Dan peta sisa buku ini.}

\framekicker{Peta}
\begin{frame}[shrink=25]{Tiga arah, dan masing-masing punya babnya sendiri}
\begin{minipage}[t]{0.3253\textwidth}
\begin{icardaccent}
\cardh{Banyak agen}
Beberapa agen dengan peran dan alat berbeda. Menarik di papan tulis, mahal di tabel biaya --- dan \textbf{bab 8} menghitungnya.
\end{icardaccent}
\end{minipage}%
\hfill
\begin{minipage}[t]{0.3253\textwidth}
\begin{icardaccent}
\cardh{Multi-modal}
Masukannya bukan cuma teks: gambar, layar, dokumen. Yang berubah bukan gelungnya, melainkan apa yang bisa jadi bukti. \textbf{Bab 9}.
\end{icardaccent}
\end{minipage}%
\hfill
\begin{minipage}[t]{0.3253\textwidth}
\begin{icardaccent}
\cardh{Agen kode}
Kode adalah alat yang paling kuat dan paling tajam: ia bisa melakukan apa pun yang bisa dilakukan kode. \textbf{Bab 10}.
\end{icardaccent}
\end{minipage}%
{\fontsize{9.4}{12.6}\selectfont\color{fg2}Ketiganya bukan jenis agen yang berbeda. Ketiganya \textbf{gelung yang sama} dengan ruang tindakan yang berbeda --- yang berubah selalu alatnya, bukan mesinnya.\par}\vskip 4pt
\end{frame}

\framekicker{Peta}
\begin{frame}[shrink=25]{Urutan babnya bukan urutan membangunnya}
{\fontsize{9.4}{12.6}\selectfont\color{fg2}Buku ini bergerak dari bagian ke sistem: model (bab 2--3), lalu yang ditambahkan padanya --- memori (4), alat (5), rencana (6) --- lalu cara menilainya (7), lalu ketiga spesialisasi (8--10).\par}\vskip 4pt
\begin{isteps}
\item \{'h': 'Membangunnya justru dari belakang', 'p': 'Kumpulan uji dulu --- sebelum agennya ada. Dua puluh kasus nyata dengan hasil yang sudah diketahui, ditulis selagi masih jujur tentang apa artinya benar.'\}
\item \{'h': 'Lalu alat, lalu batasnya', 'p': 'Alat baca dulu. Alat tulis belakangan, dan tiap satu dengan alasan tertulis.'\}
\item \{'h': 'Gelungnya paling akhir', 'p': 'Bagian yang paling menarik dibaca adalah bagian yang paling sedikit menentukan keberhasilannya.'\}
\end{isteps}
\begin{iband}[signal]
Kalau urutan membangunnya dibalik mengikuti urutan babnya, yang terjadi biasanya: gelung yang bagus, alat seadanya, dan tidak ada cara mengetahui apakah semuanya bekerja.
\end{iband}
\end{frame}

\framekicker{Penutup}
\begin{frame}[shrink=25]{Yang dibawa pulang dari bab ini}
\begin{isteps}
\item \{'h': 'Satu sifat', 'p': 'Agen memilih urutan langkahnya \textbf{saat berjalan}. Sisanya akibat.'\}
\item \{'h': 'Gelung tanpa syarat henti akan berputar', 'p': "Anggaran dan 'tidak ada alat untuk ini' adalah dua penyelamat yang paling sering lupa dipasang."\}
\item \{'h': 'Orkestrator adalah kode biasa', 'p': 'Kalau ia tipis, semua masalah akan terlihat seperti masalah model.'\}
\item \{'h': 'Otonomi itu tingkatan', 'p': 'Pisahkan alat baca dari alat tulis. Perubahan sehari, keamanan naik banyak.'\}
\item \{'h': 'Nilai dari jejaknya', 'p': 'Keluaran benar dari alasan yang salah akan berulang, dan tidak akan tertangkap kumpulan uji.'\}
\end{isteps}
\end{frame}
\note{Tutup dengan pertanyaan: dari sistem yang sedang kalian bangun, bagian mana yang benar-benar butuh memilih langkahnya sendiri? Jawabannya hampir selalu lebih kecil daripada dugaan awal.}

\end{document}
